\section{Probabilistic computation}\label{sec:comp_sampling}
The stochastic dynamical exploration of an information landscape is a powerful paradigm of general computation, and is known as probabilistic computation. Such a landscape in this context is a representation of a target probability distribution. Sampling processes represent and evaluate these distributions through the efficient production of informative stochastic samples. Uncertain, high-dimensional, and ambiguous inference problems require efficient and accurate characterisations of target distributions. Sampling methods aim to draw sampled independently from a target distribution. This paradigm has become widely utilised in machine learning and statistics, where generative models learn patterns from sparse samples of large multivariate distributions inferred from observational data \cite{delacour_neural_1997, cranmer_advances_2023}. This framework has recently gained significant traction as a basic scheme to describe biological computation. The variability in firing patterns of neural populations would thus maintain an internal model of perceptual uncertainty and evolving understanding, giving a functional role to chaotic neural activity \cite{rao_predictive_1999, koch_predicting_1999, lange_bayesian_2020}. This section outlines the paradigm of neural sampling as an interpretation of the variability in neuron firing patterns and visual search gaze trajectories (\cref{subsec:neur_sampling}) as well as their extension to describe visual searching as a sampling process (\cref{subsec:visual_sampling}).
%\subsection{Active perception}\label{subsec:active_inf}
%The dominant framework in modern computational neuroscience is that of the Bayesian brain. Karl Friston, the most cited neuroscientist \cite{Bohannon2016-kp}, developed this framework, and shifted the view of the perceptual loop of the brain. Friston's theory of active inference described perception through inferring and generating top-down sensory predictions of the world, and carrying back the bottom-up errors from that perception. This development brought a strong shift in the literature, ``marking a watershed between 20th-century thinking about the brain as a glorious stimulus–response link and more constructivist 21st century perspectives that emphasized an active sampling of the sensory world'' (\cite{friston_does_2018} pg. 1019). The active inference model of an updating internal mental state filled the explanatory gap in the understanding of the functional role of many key feedback mechanisms in the brain \cite{rao_predictive_1999, koch_predicting_1999}, and went on to yield successful predictions of how the brain optimally deciphers noisy multimodal sensory inputs \cite{kording_bayesian_2004}, as well as provide rigorous frameworks to base probabilistic computation in the brain \cite{lange_bayesian_2020}. There has recently been opposition to the general framework of the Bayesian brain and active inference \cite{mangalam_myth_2025}. Much of the arguments against such a framework arise from the philosophy of science, where active inference in perception may be supported post-hoc through ``just so'' definitions of priors \cite{mangalam_myth_2025, bowers_bayesian_2012}. It is therefore important to understand the Bayesian brain hypothesis as a normative framework from which to build falsifiable, understandable, and biologically plausible theories of neural computation.

\subsection{Neural sampling}\label{subsec:neur_sampling}
The framework of neural sampling posits this paradigm that the variability of cortical neural responses is an important functional feature of biological computation. Otherwise regarded as noise, stochastic neural activity patterns are thus presented as samples drawn from a posterior distribution over possible sensory interpretations \cite{orban_neural_2016, NIPS2002_a486cd07}. The neural sampling hypothesis has been shown to perform and predict network-level empirical patterns of neural activity in the V1, and recently by Terada and Toyoizumi in biologically realistic cortical networks, better than alternative theories \cite{NIPS2002_a486cd07, ujfalussy_sampling_2022, terada_chaotic_2024, orban_neural_2016}. Perceptual uncertainty is then encoded by the variability of these cortical responses across time and across multiple trials, with higher uncertainty reflected in wider response distributions \cite{NIPS2002_a486cd07, delacour_neural_1997, orban_neural_2016}. Neural sampling thus provides an interpretation for the way the brain navigates uncertainty at the fundamental scale of neural computation.

Studies into neural sampling show that variability is not strictly a result of explicit stochastic noise, as traditionally believed. Instead, recent work demonstrates that strongly recurrent, balanced inhibitory-excitatory deterministic neural networks can generate the high-dimensional variability which stochastic models are built to capture \cite{terada_chaotic_2024, NIPS2002_a486cd07}. Terada and Toyoizumi link this sampling of a posterior distribution to Bayesian learning through training a recurrent network to sample static features and dynamic trajectories, generalising learned priors to novel, uncertain environments through biologically plausible learning \cite{terada_chaotic_2024}. Further biologically plausible implementations of such Bayesian sampling are energy-based models (EBMs) which suggest that fast-spiking interneurons can estimate the local partition function to enable fast, localised sampling dynamics \cite{dong_neural_nodate}. The sampling hypothesis has thus given a convincing theory of biological neural computation, with testable and well-supported evidence that neural variability has an important functional role in computation \cite{orban_neural_2016, ujfalussy_sampling_2022}. Traditional Gaussian and Poisson noise models representing these stochastic network dynamics inherently correspond to continuous struggle to capture the full dynamics of complex interconnected inhibitory-excitatory (IE) biological neuronal network \cite{koren_efficient_2025, kilgore_biologically-informed_2025, song_training_2016}. Further study into the dynamics of these IE networks has thus motivated further functional models of neural activity.

\subsubsection*{Fractional neural sampling}

%In reality, neural activity has been found to exhibit structure across multiple spatial and temporal scales, often described by scale-free statistics \cite{boonstra_scale-free_2013, he_scale-free_2014, he_temporal_2010, jiao_biologically_2026, jones_scale-free_2023, lin_scale-free_2016, mashour_conscious_2020, roberts_heavy_2015, zhang_heavy-tailed_2025}. This observation, together with the neural sampling hypothesis, motivates the design and theory of Fractional Neural Sampling (FNS) \cite{qi_fractional_2022}. FNS describes the complex spatiotemporal dynamics of spiking neural circuits, where localised activity patterns propagate through superdiffusive Lévy flights. Such Lévy flight dynamics of the centre of mass of firing spikes $\textbf{x}_t$ with a Langevin equation driven by Lévy noise and coupled to a momentum term $\textbf{p}_t$, which we shall call the FNS equations of motion,
%
%\begin{equation}
%\label{eq:fns}
%\begin{aligned}
%	d\textbf{x}_t &= \gamma b(\textbf{x}_t)dt + \beta \textbf{p}_t dt + \gamma^{1/\alpha} dL_t^{\alpha}, \\
%	d\textbf{p}_t &= \beta b(\textbf{x}_t) dt
%\end{aligned}
%\end{equation}
%
%where $\textbf{x}_t$ is the (usually 2 dimensional) vector of position, $\textbf{p}_t$ is a generalised auxiliary variable analogous to a momentum of the same dimensionality, $\beta > 0$ is a parameter defining the strength of coupling between $\textbf{p}_t$ and $\textbf{x}_t$, $\gamma > 0$ is the noise strength parameter, $ 0 < \alpha \leq 2$ is the Lévy tail exponent, and $dL_t^{\alpha}$ is the Lévy stable noise drawn from a $\mathcal{S}\alpha\mathcal{S}$ distribution. The term $b(\textbf{x}_t)$ captures a deterministic gradient term in the dynamics, defined such that the invariant distribution of $\textbf{x}_t$ after long time $t$ approaches the target distribution sampled, $\pi(\textbf{x})$. The exact expression for $b_{\alpha}(\textbf{x})$ is derived in \cite{qi_fractional_2022}, and approaches the gradient of the logarithm of the density as $\alpha \rightarrow 2$, with the full form for any $ 0 < \alpha \leq 2$ given in terms of the $\alpha-2$-order Riesz fractional derivative $\mathcal{D}^{\alpha - 2}_{\textbf{x}}$ as
%
%\begin{equation}\label{eq:full_drift}
%	b_{\alpha}(\textbf{x}) = \frac{\mathcal{D}^{\alpha - 2}_{\textbf{x}}[\pi(\textbf{x}) \partial_{\textbf{x}} \ln\pi(\textbf{x})]}{\pi(\textbf{x})}
%\end{equation}
%with the useful local approximation, which becomes exact at $\alpha=2$,
%\begin{equation}\label{eq:drift_approx}
%    b_{\alpha}(\textbf{x}) \approx c_{\alpha} \nabla_{\textbf{x}} \ln \pi(\textbf{x}), \qquad c_{\alpha} = \frac{\Gamma(\alpha - 1)}{\Gamma(\frac{\alpha}{2})^2}
%\end{equation}

The FNS dynamics are called fractional, as the non-local dynamics of their distribution require fractional calculus to describe such spatial superdiffusion in Lévy flights \cite{tarasov_general_2021}. The theory of FNS thus unifies the scale-free nonlocal dynamics of neural firing patterns with the neural sampling framework derived from the ideal Bayesian scheme of active inference. The critical advantage of FNS over other standard continuous random walks is its ability to efficiently sample complex, multimodal probability distributions with far-apart modes—likely a major feature of processed natural environments replete with multiple distinct salient features. While Brownian-motion-based samplers typically remain trapped in local modes due to high potential energy barriers, the heavy-tailed, large jumps inherent in FNS enable the neural sampler to adaptively and freely switch between multiple modes \cite{neal_sampling_1996, latuszynski_mcmc_2025, simsekli_fractional_nodate, metzler_random_2000}. Furthermore, the generality of the FNS description permits the application of FNS equations of motion to gaze dynamics, as the motion of the eyes are driven by an IE neural circuit which is expected to keep key features of the FNS model invariant, such as superdiffusion, heavy-tailed statistics, and nonlocal interactions.

\subsection{Visual sampling}\label{subsec:visual_sampling}
This thesis presents visual search tasks as such a paradigm where a sampling model is well motivated, and driven by the underlying neural sampling process as well. The visual system is evolved to actively and efficiently process noisy information and evidence from the external world, and plan future search paths to gather further noisy evidence \cite{Renninger2005-ee, renninger_where_2007, najemnik_optimal_2005, bujia_modeling_2022}. In the context of visual searching, visually interesting targets with \textit{saliency} as the equivalent of the target landscape, and where the information is located in a target image. Here, saliency is equivalent to the Boltzmann log-likelihood potential from the probability distribution of fixations $p(\textbf{x})$
\begin{equation}\label{eq:saliency}
	S(\textbf{x}) = -\log(p(\textbf{x}))
\end{equation}
where visual search paths are most likely to be drawn to is a common technique whic \cite{Bruce_2009_Saliency}. As such, undirected visual searching of an image will show the areas most salient to the brain, allowing insight into the computational sampling mechanism of the visual system. In particular, search tasks require sophisticated integration of the cost and expected information gain of decisions, which invites fruitful decision analysis \cite{Wedel_2023}. Empirical experiments show that eye movements show good agreement with greedy local information maximising fixations, given well-motivated visual sensitivity based on human physiology \cite{Radulescu_2026, najemnik_optimal_2005, renninger_where_2007}.

%%This dynamical switching provides a functional explanation for the bimodality of estimate distributions observed in human perceptual switching and visual perception inference.
%\subsection{Markov Chain Monte Carlo methods}\label{subsec:MCMC}
%This view of neural- and visual sampling of a target distribution motivates analogy to a foundational set of methods in machine learning. Markov Chain Monte Carlo (MCMC) methods have become the algorithmic basis for generating samples from target distributions in rigorous Bayesian inference in machine learning \cite{cranmer_advances_2023, latuszynski_mcmc_2025, speagle_conceptual_2020, andrieu_particle_2010}. MCMC operates by constructing a Markov chain, where the transition matrix is calculated such that the equilibrium distribution of samples is the desired target distribution. In the context of neural computation, biological circuits are often modelled as implementing continuous-time variants of MCMC sampling of neural population targets, most notably through Langevin sampling and Hamiltonian Monte Carlo (HMC) \cite{latuszynski_mcmc_2025, betancourt_conceptual_2018}. Langevin sampling dynamics utilise local gradient information from the probability landscape, combined with Gaussian white noise, to explore the state space over time. Such dynamics are identical to the Langevin system studied in statistical physics, described in \cref{eq:langevin}, where $V(\textbf{x}) = - \log(\pi(\textbf{x}))$ is the potential function over parameters $\textbf{x}$ required yielding a distribution $\pi(\textbf{x})$ of samples. Despite their ubiquity, traditional MCMC methods encounter severe theoretical and practical limitations when exploring unknown, heavy-tailed, and/or multimodal target distributions. Standard Langevin and HMC dynamics are driven by continuous Brownian and Hamiltonian motion respectively, and the traversal of low-probability regions, equivalent to dynamical energy barriers, becomes exponentially slow. Consequently, the mean exit time for the sampler to leave one mode and discover another grows exponentially with the modal separation distance \cite{latuszynski_mcmc_2025}. To accelerate this sluggish exploration, algorithms often introduce auxiliary momentum variables to drive global landscape traversal, as in the FNS equations of motion \ref{eq:fns}. However, even with momentum-based acceleration, traversing deep energy valleys remains highly inefficient \cite{latuszynski_mcmc_2025, neal_sampling_1996, betancourt_conceptual_2018}. The nonlocal dynamics introduced in fractional sampling are one branch of research into how rapid mixing of samples in complex distributions can be achieved without the need for nontrivial parameter tuning and knowledge of the distribution of interest. For instance, fractional algorithms based on Hamiltonian Monte Carlo sampling and second-order Langevin sampling with a momentum variable have been briefly developed and explored \cite{simsekli_fractional_nodate, ijcai2018p419, qi_fractional_2022, oechsler_levy_2024, behme_levy_2026}. Importantly, FNS dynamics over a probability landscape yield a mean exit time between modes which scales linearly rather than exponentially with separation \cite{qi_fractional_2022}. FNS dynamics thus demonstrate a profound computational advantage for efficient biological inference through it's background as an effective sampling algorithm.